%==========================================================
\section{Teaching the AI with more varied weather}
\label{sec:varied}
%==========================================================

Two questions were left open above. Why does the AI still measure too little rain after a rainy day in Italy
(Section~\ref{sec:drylevel})? And why do training runs that agree on \emph{whether} it rains disagree so much on
\emph{how much} (Section~\ref{sec:se_results})? We first looked for the causes, then tried fixes, and kept the test months
untouched throughout.

\subsection{Why the training runs disagree on amounts}
\label{sec:why_runs}
A run's rain total can be wrong in three ways: it calls rain ``dry'' (missed rain), it calls a dry moment ``rain'' (false
rain), or it says ``rain'' at the right moments but with the wrong strength. We split the total of each of the five runs
of Section~\ref{sec:se_results} into these three parts, against the radar, on the July validation weeks and on the
August test month:
\begin{center}\footnotesize\setlength{\tabcolsep}{5pt}
\begin{tabular}{cccccc}
\toprule
& \multicolumn{2}{c}{\textbf{July (validation)}} & \multicolumn{3}{c}{\textbf{August (test)}}\\
\cmidrule(lr){2-3}\cmidrule(lr){4-6}
Run & bias & strength when it says rain & bias & strength & missed / false rain\\
\midrule
\SeedRows
\bottomrule
\end{tabular}
\end{center}
(``Strength'' is the run's rain divided by the radar's, at the moments when both say it rains; 1.00 is perfect. The
missed and false columns are shares of the radar's total.)

Three things stand out:
\begin{itemize}
\item \textbf{It is the strength, not the switch.} All runs miss about the same rain (10--12\%) and add about the same
false rain. What differs is how strongly they convert a decibel into millimetres: from 0.69 to 0.89 of the radar in August.
Fixing the rain/no-rain threshold at the same value for every run does not change this.
\item \textbf{July tells us about August.} A run that reads high in July also reads high in August: the correlation between
the two columns of bias is \SeedValTestR. So the spread is a property of each trained network, not bad luck in August.
\item \textbf{All runs read lower in August than in July}, by 15 to 20 points. August 2015 had more heavy downpours, and the
strength is learned mostly from the few heavy minutes of six weeks in one city. Each run fits those few minutes a little
differently.
\end{itemize}
That last point is the cause: too little and too uniform training rain. The cure we expected was more varied rain.

\subsection{What a long rain looks like to a link}
\label{sec:may6}
The Italian ``after a mostly rainy day'' hours are not spread over the summer. In 2022, 85\% of them come from a single
storm: a day and a half of steady rain on 5--7 May over the Apennine hills. Figure~\ref{fig:may6} shows one link through it,
next to a rain gauge 508\,m away.

\begin{figure}[h]
\centering
\begin{tikzpicture}
\begin{groupplot}[group style={group size=1 by 2,vertical sep=1.5cm,xlabels at=edge bottom},
  width=0.95\textwidth,height=4.6cm,xmin=0,xmax=60,grid=major,grid style={inkc!12},
  xtick={0,12,24,36,48,60},xticklabels={{5 May\\12:00},{6 May\\00:00},{6 May\\12:00},{7 May\\00:00},{7 May\\12:00},{8 May\\00:00}},
  xticklabel style={align=center,font=\scriptsize},
  tick label style={font=\small},label style={font=\small},
  legend style={font=\scriptsize,draw=none,fill=white,fill opacity=0.85,text opacity=1},legend cell align=left]
\nextgroupplot[ylabel={rain (mm/h)},ymin=0,ymax=4.5,legend pos=north east]
  \addplot[thick,color=gaugec,mark=*,mark size=1pt] table[x=hour,y=gauge] {may6.dat};
  \addlegendentry{rain gauge}
  \addplot[thick,color=radarc] table[x=hour,y=radar] {may6.dat};
  \addlegendentry{radar along the link}
  \addplot[thick,color=inkc,dashed] table[x=hour,y=ai_se] {may6.dat};
  \addlegendentry{AI, Sweden only}
  \addplot[thick,color=rainc] table[x=hour,y=ai_cv] {may6.dat};
  \addlegendentry{AI, varied training (never saw May)}
\nextgroupplot[ylabel={dB},xlabel={hourly averages (UTC)},ymin=0,ymax=6,legend pos=north east]
  \addplot[thick,color=inkc] table[x=hour,y=excess] {may6.dat};
  \addlegendentry{extra loss above the dry level}
  \addplot[thick,color=radarc,dashed] table[x=hour,y=expected] {may6.dat};
  \addlegendentry{loss the radar's rain should cause}
\end{groupplot}
\end{tikzpicture}
\caption{The long storm of 5--7 May 2022 on sub-link 182\_1 (5.3\,km, 24.6\,GHz). Top: rain. Bottom: what the link
measured. Storm totals: gauge \Maygauge\,mm, radar \Mayradar\,mm, AI trained on Sweden only \Mayaise\,mm, AI with varied
training \Mayaicv\,mm.}
\label{fig:may6}
\end{figure}

What the figure, and the same numbers for every link, show:
\begin{enumerate}
\item \textbf{The dry level is not the problem.} In every group of Italian rainy hours, the 3-day low value is at or a little below (by
0.1--0.2\,dB, and up to 0.9\,dB in the long storm) the level the link has when the radar says it is dry. It is never too
high, so it cannot be what makes the AI read too little.
\item \textbf{The link sees more than the rain.} While it rains, the extra loss is the rain's share plus a steady 1.5--2\,dB
more: water on the antenna covers, which stay wet in a long rain (the wet antenna of Section~\ref{sec:fakes}).
\item \textbf{The extra loss lasts a whole day after the radar stops} (6 May evening to 7 May noon). The gauge shows why:
light rain, 0.2--1.2\,mm/h, kept falling. The radar, looking from far above, saw almost none of it.
\item \textbf{The AI gives up.} After about twelve hours of raised loss, the AI reports almost nothing, even at hours when
the radar shows 2--4\,mm/h. To a single link, ``my loss has been raised for half a day'' looks like a faulty or drifting link, and
our training data are full of those. That is what the AI learned.
\end{enumerate}
So the cause is not one bug. Three things push the same way: a long rain wets the antennas, the radar we train on misses the
light rain at its end (so the training answer says ``dry'' while the link is still wet), and faulty links teach that a
loss which stays raised for hours is not rain.

\subsection{Trying fixes, one Italian summer at a time}
Each fix was trained five times (seeds 0--4) on Sweden June--July, plus Italy 2021 where stated, and tested on what it never
saw: Sweden in August and Italy in 2022. ``Rain after a mostly rainy day'' is measured here in 15-minute steps, which is
stricter than the hourly table of Section~\ref{sec:drylevel}.
\begin{center}\footnotesize\setlength{\tabcolsep}{3.5pt}
\begin{tabular}{p{5.1cm}ccccc}
\toprule
& \multicolumn{2}{c}{\textbf{Sweden, August}} & \multicolumn{3}{c}{\textbf{Italy 2022}}\\
\cmidrule(lr){2-3}\cmidrule(lr){4-6}
& & bias vs gauges, & total bias, & rain after a & rain after a\\
Training & MCC & 5 runs & 5 runs & dry day & mostly rainy day\\
\midrule
\ExpRows
\bottomrule
\end{tabular}
\end{center}
\begin{itemize}
\item \textbf{One Italian summer fixes most of the spread.} Five runs that disagreed by 47 points on Italy's 2022 total
now agree within 7 to 12 points, and in Sweden within 5 to 17 points instead of 26. Rain detection in Sweden improves too
(MCC 0.65 against 0.62), even though no Swedish data was added.
\item \textbf{Training only on reliable links helps the long storm.} With the reliability score of
Section~\ref{sec:reliability} applied to the Italian training links, the share of the long storm measured rises from 28\% to
42\%, with nothing else lost. This supports point 3 above: faulty links were teaching the AI to ignore long raised loss.
\item \textbf{Two ideas did not work.} Scoring the rate error in a log scale (so light rain counts as much as heavy rain)
made the AI read 5--15\% too little in Italy. Giving it two extra inputs from its neighbours (``how many links around me
are raised right now?'') improved rain detection (MCC 0.68) but cost amounts elsewhere and barely helped the long storm.
\end{itemize}

\subsection{The final check: every Italian month, judged by a model that never saw it}
\label{sec:cv}
One Italian summer may be luck. So we used both summers, month by month: to judge May, we train on Sweden plus June to
September of \emph{both} Italian years; to judge June, on the other four months; and so on. This is called
\textbf{cross-validation}. Every Italian month of 2021 and 2022 is then judged by a network that never saw it, and each
network learns from long rains of the other months (for example September 2021 when May is judged). Training uses the
reliable links only; the last week of each training month is kept apart to pick the checkpoint and the threshold; three
runs per month. Sweden's August is never used.

The same hourly table as in Section~\ref{sec:drylevel}, same links, same groups (rain measured as a share of the radar's):
\begin{center}\footnotesize
\begin{tabular}{lcccc}
\toprule
& \multicolumn{4}{c}{\textbf{share of the previous 24 hours that was rainy}}\\
\cmidrule(lr){2-5}
& under 10\% & 10--25\% & 25--50\% & over 50\%\\
\midrule
Italy 2021, Sweden-only AI & \SeOnlyLongAa & \SeOnlyLongAb & \SeOnlyLongAc & \SeOnlyLongAd\\
Italy 2021, \textbf{varied training} & \textbf{\CVrelLongAa} & \textbf{\CVrelLongAb} & \textbf{\CVrelLongAc} & \textbf{\CVrelLongAd}\\
\midrule
Italy 2022, Sweden-only AI & \SeOnlyLongBa & \SeOnlyLongBb & \SeOnlyLongBc & \SeOnlyLongBd\\
Italy 2022, \textbf{varied training} & \textbf{\CVrelLongBa} & \textbf{\CVrelLongBb} & \textbf{\CVrelLongBc} & \CVrelLongBd\\
\bottomrule
\end{tabular}
\end{center}
Against the rain gauges, over all rainy hours, the varied AI measures \CVrelGaugesA{} (2021) and \CVrelGaugesB{} (2022) of
the gauges' rain. The three runs per month agree on Italy's total within \CVrelBiasA{} (2021) and \CVrelBiasB{} (2022).
In Sweden's August, the same networks score MCC \CVrelSeMCC, hourly correlation with gauges \CVrelSeR, error
\CVrelSeRMSE\,mm/h and bias \CVrelSeBias{} (runs from \CVrelSeBiasRange), against 0.615, 0.79, 1.04\,mm/h and $-20\%$ for the
Sweden-only model: rain detection and agreement between runs improve clearly, the Swedish shortfall only a little.

\begin{careful}
\textbf{What this test is, and is not.} The Italian links are the same in every month: this tests new weather on known links,
not new links. The Italy and Netherlands results of Sections~\ref{sec:italy}--\ref{sec:neighbours}, and the maps of Part~3,
keep the Sweden-only model on purpose, so that they remain tests on countries the AI has never seen. For Algeria, which
has no reference data of its own, the recipe of this section is the one to use: train on several countries, reliable links
only, and average several runs.
\end{careful}

\begin{oneline}
More varied rain fixed most of both problems: training runs now agree on amounts, Sweden improved, and in Italy the AI
measures \CVrelLongAd{} of the rain after a rainy day in 2021. Two things remain: the one long storm of May 2022 is only
partly measured (\CVrelLongBd{}), and in Sweden's August the AI still reads \CVrelSeBias{} low.
\end{oneline}
